\documentclass[aspectratio=169]{beamer} % 16:9 is standard for modern projectors
\usetheme{moloch} % or \usetheme{metropolis}

% Optional: Fira Sans font (requires XeLaTeX or LuaLaTeX)
\usepackage[sfdefault]{FiraSans}
\usepackage{graphicx}

% Customizing colors for a "systems" feel (Deep Slate / Cyan)
\definecolor{SystemSlate}{HTML}{23373B}
\definecolor{HighlightCyan}{HTML}{00B4D8}

\setbeamercolor{frametitle}{bg=SystemSlate, fg=white}
\setbeamercolor{progress bar}{fg=HighlightCyan, bg=SystemSlate}

\title{Overheads of Hardware vs. Software Virtualization}
\subtitle{A quantitive analysis of KVM and QEMU TCG Memory subsystems}
\author{Milad Zarei - Taha Biklariyan - Nazanin Sharifi}
\date{Semester 4042}
\institute{OS Tracing - Dr. Azhari @ IUST}

\begin{document}
  \maketitle
  \begin{frame}{Background}
      \begin{columns}[T] % [T] aligns columns at the top
          
          % Left Column: Content
          \begin{column}{0.6\textwidth}
              \begin{itemize}
                  \item Virtualization is widely-used in cloud computing, OS development, etc.
                  \item Allows you to run multiple \textit{guest operating systems} inside a \textit{host operating system}.
                  \item \textbf{Hardware-Assisted Virtualization:}
                  \begin{itemize}
                      \item Powered by KVM in the Linux kernel for virtualization on the same host architecture.
                      \item Processing is done on actual hardware via CPU Virtualization technologies (VT-x/AMD-V).
                  \end{itemize}
              \end{itemize}
          \end{column}
          
          % Right Column: Figure
          \begin{column}{0.4\textwidth}
              \centering
              % Replace 'virt-stack' with your actual image filename
              \includegraphics[width=\textwidth]{figures/virtualization.webp}
              \vspace{0.5cm} % Add some breathing room
              \small
              \textit{Figure 1: KVM Virtualization Stack Architecture}
          \end{column}
          
      \end{columns}
  \end{frame}

  \begin{frame}{Background}
    \begin{itemize}
      \item Virtualization is widely-used in cloud computing, OS development. etc.
      \item Allows you to run multiple \textit{guest operating systems} inside a \textit{host operating system}.
      
    \end{itemize}
  \end{frame}
  \begin{frame}
    \begin{itemize}
      \item Most common methods of virtualization:
        \begin{itemize}
          \item \textbf{Hardware-Assisted}: Powered by KVM in the Linux kernel for virtualization that runs on the same architecture as the host.
            Most of the processing is done on the actual hardware, alongside the host! By the help of CPU Virtualization technologies.
          \item \textbf{Emulation:} For running OSes that run on entirely different CPU architectures! Powered by QEMU's TCG (Tiny Code Generator). 
        \end{itemize}
    \end{itemize}
  \end{frame}
  \begin{frame}{Motivation \& Problem}
    introduce the challenge of virtual memory and two approaches of ept and softmmu
    introduce the challenge of
    translation of instructions and how tcg is different from kvm. (Not very complex)
  \end{frame}
  \begin{frame}{Goal}
    Explain the goal of trying to measure differences via tools like perf and benchmarks
  \end{frame}
  \begin{frame}{System Setup}
    Explain the cpu isolation (hyper threading disabled), running debian nocloud images via scripts and communicating with them via ssh and perf. one is arm and one is x86
  \end{frame}
  \begin{frame}{Benchmarks}
    Explain the benchmarks that we have with brief explanation and what stress do they put.
  \end{frame}
  \begin{frame}{Benchmark 1 results and analysis}
  \end{frame}
  \begin{frame}{Benchmark 2 results and analysis}
  \end{frame}
  \begin{frame}{Benchmark 3 results and analysis}
  \end{frame}
  \begin{frame}{Benchmark 4 results and analysis}
  \end{frame}
\end{document}
